% =====================================================================
% DRAFT MANUSCRIPT — Nature Neuroscience-style
% Generated as a writing scaffold (2026-06-24). All paragraphs drafted;
% numerics are real (sourced from the headline files / ledgers / audit CSVs).
% \TODO{...} marks a value NOT in the crystallization context — fill from the
% cited CSV before submission; do NOT invent.
%
% Provenance for every number is in a % comment next to it (audit_NN / CSV).
% Rules honoured: results stated in the LRG (cophenetic/Grassmann) framework,
% raw FC is the baseline; matched-strength (C3) is the referee for every cohort
% claim; OFC is a CONCENTRATION not a container; no behavioural claims (no data);
% inference->cingulate is a directional hint only.
% =====================================================================
\documentclass[11pt]{article}

\usepackage[margin=1in]{geometry}
\usepackage{amsmath,amssymb}
\usepackage{booktabs}
\usepackage{graphicx}
\usepackage{xcolor}
\usepackage{lineno}
\usepackage[hidelinks]{hyperref}
\linenumbers

% --- convenience macros ---
\newcommand{\rhocoph}{\ensuremath{\rho^{\mathrm{coph}}}}
\newcommand{\dG}{\ensuremath{d_{G}}}
\newcommand{\rhat}{\ensuremath{\hat{\rho}(\tau)}}
\newcommand{\TODO}[1]{\textcolor{red}{[\textbf{TODO}: #1]}}
\newcommand{\fig}[1]{Fig.~\ref{#1}}
\newcommand{\tab}[1]{Table~\ref{#1}}

\title{Offline abstraction of a learned structure leaves a multiscale,
band-specific trace in human cortical connectivity}
% alt title: The resting human brain holds an offline, multiscale trace of what it inferred
\author{\TODO{author list}}
\date{Draft --- \today}

\begin{document}
\maketitle

% =====================================================================
\begin{abstract}
\noindent
The brain must consolidate not only what it experiences but what it
\emph{infers}. Whether inferred relational structure leaves an offline trace in
human cortical connectivity --- and at what organisational scale --- is unknown,
in part because connectivity is usually read at a single threshold or through
global spectral embeddings that are scale-blind and hard to interpret. We
recorded stereo-EEG in ten patients across the four phases of a transitive-
inference task (rest $\rightarrow$ learning $\rightarrow$ test $\rightarrow$ rest)
and read functional connectivity through the full multiscale hierarchy of a
graph-Laplacian diffusion operator, with two non-redundant probes. We find,
first, that the task leaves a band-specific ($\beta$) multiscale reorganisation
that persists into post-task rest, survives a strength-matched surrogate null,
concentrates in --- but is not confined to --- orbitofrontal cortex, is carried
by healthy gray-matter coupling, and is \emph{held} as a sustained, non-ergodic
state (10/10 patients, $p=0.001$) rather than re-emerging in transient bursts.
Second, this persistent trace carries an inference-specific component beyond mere
re-exposure to the premises ($\beta$ only, $p\approx0.007$), favouring the
mesoscale of multi-step integration and anchored in an orbitofrontal map that
learning itself establishes --- an offline signature of \emph{abstracting} a
learned order. Third, the same diffusion operator, read as inter-contact
affinity, localises epileptogenic tissue (leave-one-patient-out
AUC~$\approx0.81$, precision@5~$\approx60\%$). The reorganisation is invisible to
edge-wise and to spectral-clustering/PCA analyses, which a multiscale read-out
recovers. Cognition and clinic thus emerge from one operator.
\end{abstract}

% =====================================================================
\section{Introduction}

Learning a transitive order (A$>$B, B$>$C, $\dots$) lets an animal answer a pair
it has never seen (A vs.\ D). This is not memory for experience but the
construction and use of an abstract relational structure --- a ``cognitive
map'' of task space. A central prediction of consolidation theory is that such
structure, once built, is rehearsed and stabilised offline. In humans this has
been inaccessible at the network level: scalp recordings lack the spatial
resolution, and intracranial connectivity has been read with tools that collapse
the brain's organisation onto a single scale.

Functional connectivity is most often summarised by thresholding an adjacency
matrix at one density, or by embedding its leading eigenmodes (spectral
clustering, PCA). Both choices discard the fact that brain networks are organised
\emph{across} scales --- from tightly coupled local cliques to broad
inter-system communities --- and both are hard to interpret anatomically. A
reorganisation that lives in the relationship \emph{between} scales, rather than
at any one of them, is therefore easy to miss.

We read connectivity differently. From a single graph-Laplacian diffusion
operator, $\rhat = e^{-\tau\hat{L}}/Z$, we obtain a continuous family of
representations indexed by diffusion time $\tau$ --- a multiscale renormalisation
of the network --- and summarise it with two complementary probes: a per-pair
\emph{cophenetic} distance (\rhocoph), which is multiscale and locally
interpretable, and a global-mode \emph{Grassmann} distance (\dG), which is the
subspace view that ordinary spectral methods approximate. We use a
volume-conduction-immune connectivity substrate (the band-averaged magnitude of
the imaginary part of coherency) and a locked battery of null models, of which a
strength-matched surrogate is the binding referee.

Recording across rest, learning (premise pairs shown --- encoding), test (novel
pairs reasoned out --- inference), and post-task rest, we ask three questions.
Does the task leave a trace that persists into rest, and what is its scale,
spectral signature, anatomy, and dynamics? Does that trace carry the
\emph{inferred} structure specifically, beyond what was merely seen? And does the
same operator, applied to the same recordings, speak to the clinical problem of
localising the seizure-onset zone?

% =====================================================================
\section{Results}

% ----- cohort sentence (shared) -----
All analyses are over the full cohort of $n=10$ patients
(\TODO{confirm IDs: Pat\_02, 03, 05, 06, 07, 08, 10, 13, 14, 15}). Results are
stated in the Laplacian (cophenetic / Grassmann) framework; raw connectivity
serves only as a baseline. Every cohort-level claim is evaluated against a
strength-matched surrogate that preserves each node's total connectivity while
destroying higher-order structure (control C3; see Methods); within-baseline
split-half (C1), drift (C2), cross-probe (C4) and epileptic-zone-exclusion (C5)
controls accompany it.

% =====================================================================
\subsection{A held, multiscale trace of learning and reasoning}

\paragraph{Reasoning leaves an imprint that rest does not erase.}
Comparing the resting network before and after the task, the post-task resting
configuration sat reliably closer to the task configuration than the pre-task
resting configuration did, at the cohort level. On the per-pair cophenetic probe
--- our primary read-out --- the effect was carried most clearly by the $\beta$
band (cohort-median $\rho_{\text{split}}=+0.22$; 7 of 10 patients above their own
strength-matched null; paired $p=0.005$), with a real but weaker $\alpha$ component
($p=0.002$) and a borderline low-$\gamma$ trace that nonetheless localises (below);
$\theta$ showed nothing.
% source: 01_trace.md N1.1; matched_strength_surrogate_split_baseline (audit_63); per_band_taxonomy_verdict (audit_83b)
The persistence is a \emph{cohort-level} effect over a genuinely variable cohort:
the between-patient spread is real, not a measurement-reliability artifact --- a
residual between-patient axis survives correction for the split-half reliability of
all four phase estimates (Methods) --- so we claim a cohort-level trace, never that
every patient expresses it.
% source: multiphase_snr (audit_145); 2026-06-25_multiphase-snr-reliability

\paragraph{The trace consolidates in the orbitofrontal cognitive map.}
Where the persistent reorganisation lands is sharply rhythm-specific and
biologically pointed. The distributed $\beta$ trace is over-expressed, above each
patient's own brain-wide baseline, in the orbitofrontal cortex (OFC) system ---
the hub for building relational ``cognitive maps'', exactly what a
transitive-inference task should leave behind. Against a matched-strength null with
$R=1000$ surrogates, OFC cleared Benjamini--Hochberg correction across nine
a-priori cortical systems ($q\approx0.009$--$0.013$) in all four analysis
conditions (contact- vs.\ shaft-level $\times$ epileptic-zone included/excluded),
was leave-one-out immovable (top-ranked in 10/10 folds) and shaft-collapse robust,
and sat in low-strength, non-hub, bilateral tissue --- with primary sensorimotor
cortex the one depleted complement.
% source: N1.3; localization_atlas/matched_strength_R1000 + carrier_loo (audit_83/83b/92); ANATOMY_LEDGER
Low-$\gamma$ concentrated, independently, in prefrontal cortex (significant in
every leave-one-out fold, top-ranked in 8/10 --- corroborating, a rung below
$\beta$), whereas the $\alpha$ trace, though real, had no surviving anatomical home
and localised differently in each patient. The concentration is therefore a
\emph{hotspot on a distributed trace, not a container} (within-OFC lean
$\approx0.59$): it peaks at the system scale and washes out at lobe
($p\approx0.29$) and hemisphere ($p\approx0.48$--$0.76$). We frame this as
concentration above baseline, not as ``the trace lives in OFC''.
% source: carrier_loo (audit_83b, low_gamma PFC LOO); per_band_taxonomy_verdict

\paragraph{An independent read-out confirms the \texorpdfstring{$\beta$}{beta}
trace.}
The $\beta$ finding does not hinge on a single analysis. It clears the
strength-matched surrogate (control C3), so it is not a rearrangement of the
strongest connections; and a second, independent read-out of the same hierarchy
--- a global-mode subspace (Grassmann) distance --- confirms the trace at $\beta$
(cluster-extent permutation, leave-one-out robust), the convergence of two views
that are otherwise complementary across the remaining bands (\tab{tab:bandprobe};
full read-out $\times$ band map and construction in Methods).
% source: grassmann_cluster_extent (audit_70); ctm_triangle C1--C4 (audit_33/71)
The effect is genuinely multiscale: at $\beta$ the cophenetic hierarchy clears its
strength-matched control ($p=0.005$) where the matched raw-edge comparison falls
just short ($p=0.053$), and only the hierarchy expresses the absent-$\theta$ /
present-$\beta$ band structure that the raw edges flatten onto a uniform positive
floor.
% source: raw_vs_multiscale (audit_146); 2026-06-25_raw-vs-multiscale-trace

\paragraph{The \texorpdfstring{$\beta$}{beta} trace rides healthy cortex and
spares the seizure core.}
Stated in the Laplacian framework, the $\beta$ trace is carried by coupling
between healthy gray-matter contacts: gray--gray pairs form the significant
pair-count hotspot ($p<0.001$) while white-matter--white-matter pairs are
depleted, and pairs internal to the clinically labelled seizure core do not carry
it (their matched-strength contrast is non-significant, $p_{\text{pair}}\approx
0.945$; ``core spared'' is directional only)
% source: N1.4; epi_stratified/{cophenetic_cohort,pairclass_decimation_cohort}.csv (audit_77/85)
. Excluding the seizure-onset contacts altogether left the $\beta$ and $\alpha$
traces intact, so the result is neither a pathology artifact nor an
everywhere-and-nowhere effect. We explicitly withdraw two earlier sub-claims that
did not survive control: that epi-exclusion \emph{strengthens} the trace (a
generic node-count effect, ruled out by a size-matched decimation control) and
that the epi--healthy interface is its strongest carrier (it is not a pair-count
hotspot).

\paragraph{The resting brain dwells in the task's state rather than flashing back
to it.}
Resolving post-task rest into sliding windows, the per-window similarity to the
task configuration (relative to the pre-task baseline) was elevated throughout
rest, not on average alone but window after window: the contrast was positive in
10/10 patients ($p=0.001$, leave-one-out $p=0.002$)
% source: N1.6; replay_states/sustained_reinstatement_cohort.csv (audit_125/129)
. This ``hold'' was scale-invariant and representation-invariant --- positive at
every diffusion scale $\tau$ and across five distinct Laplacian read-outs
(cophenetic, magnetic/directional, raw propagator, subspace/Grassmann,
normalised; per-representation shift $+8$ to $+10$ of 10, $p\le0.007$)
% source: replay_states/{tau_resolved,directional,raw_propagator,subspace,normlap}_cohort.csv (audit_133-139)
. The complementary, transient reading was a rigorous cohort negative: across
every collapsed axis, every band, and the whole resolvable timescale from $20$~s
down to the high-$\gamma$ coherence floor near $0.2$~s, there were no bursts, no
separable states and no sequences, each tested against a time-shuffled,
node-permuted placebo null
% source: replay_states/subsecond_gamma_cohort.csv (audit_124-141); report 2026-06-22_replay-states-verification
. The post-task resting network is thus held in a sustained, non-ergodic
reorganisation --- it does not return to where it started --- rather than
re-entering the task state in brief flashes. (We claim ``held'', not ``a tighter
attractor'': the latter was only a weak trend, $p=0.08$.)

\paragraph{Over the seizure core, the memory and inference bands part ways.}
The way the persistence treats diseased tissue is itself band-dependent, and it
is the natural bridge to the clinical read-out below. In $\beta$ --- the focal,
inference-carrying band --- the seizure core is spared (above). In $\alpha$ ---
the diffuse memory band --- the same per-pair persistence instead \emph{recruits}
the epileptic core: coupling internal to the seizure zone is concentrated and
clears the matched-strength null (control C5)
% source: N1.7 bridge; verify_beta_spares_alpha_recruits.md
. The cognitive trace and the epileptogenic network are therefore not
independent, and they relate oppositely depending on rhythm.

% ---- Table 1 ----
\begin{table}[t]
\centering
\caption{\textbf{The persistence trace across frequency bands and read-outs.}
A check marks a cohort-level trace surviving the strength-matched surrogate (C3);
a parenthesised check marks a borderline cohort gate that nonetheless yields a
consistent anatomical localisation (low-$\gamma\to$PFC). The per-pair cophenetic
and global-mode Grassmann read-outs converge at $\beta$ (the flagship band) and are
otherwise complementary; this is a supporting cross-read-out summary (per-band
statistics in \TODO{supp.\ table from VERDICT\_LEDGER.md}). The Grassmann arm uses
the triangle construction and is not yet baseline-matched to the cophenetic
split-baseline (Methods).}
\label{tab:bandprobe}
\begin{tabular}{lcc}
\toprule
Band & Cophenetic \rhocoph{} (per-pair, multiscale) & Grassmann \dG{} (global modes) \\
\midrule
$\delta$       & --          & \checkmark \\
$\theta$       & --          & --          \\
$\alpha$       & \checkmark  & --          \\
$\beta$        & \checkmark  & \checkmark \\
low-$\gamma$   & ($\checkmark$) & \checkmark \\
high-$\gamma$  & --          & --          \\
\bottomrule
\end{tabular}
\end{table}

% =====================================================================
\subsection{Offline abstraction of a learned structure}
% flagship. full heading alt: "...: the trace carries the relations the brain
% inferred, not just the pairs it saw."

The persistence trace (above) establishes that the resting brain \emph{holds} a
reorganisation and that our method can see it. We next asked \emph{what} it
holds. The four-phase design is what makes this answerable: the learning phase,
in which premises are shown, is an encoding reference that can be subtracted from
the test phase, in which novel pairs must be reasoned out. Writing per-pair
cophenetic distances as $e = D_{\text{learn}} - D_{\text{pre}}$ (encoding),
$f = D_{\text{test}} - D_{\text{learn}}$ (inference-specific) and
$p = D_{\text{post}} - D_{\text{pre}}$ (persistence), the consolidation arc tests
whether task-induced change persists: an encoding echo $T_{\text{learn}}=
\mathrm{corr}(e,p)$ and an inference-specific term $T_{\text{inf}\cdot e}=
\mathrm{pcorr}(f,p\mid e)$ that partials encoding out. The decomposition
reproduces the locked persistence trace to numerical precision
($\sim\!10^{-16}$ per patient)
% source: N2.1; consolidation_arc/arc_per_patient.csv (audit_103)
.

\paragraph{The persistent trace carries what was inferred, not merely what was
seen.}
The inference-specific component itself persisted into post-task rest --- in
$\beta$, and in no other band ($p\approx0.007$, leave-one-out robust, clearing
the mandatory matched-strength null)
% source: N2.2; consolidation_arc/arc_null_per_patient.csv (audit_103 --null)
. The offline trace is therefore not merely a re-exposure echo of the premises:
it carries a component tied specifically to the inferential computation. This is
the paper's central claim --- an offline signature of abstracting a learned order
--- and we are deliberate about its ceiling. It is a signature of
inference-specific \emph{reorganisation that persists}; because task performance
was not recorded and cannot be obtained, the claim can never be tied to
inference \emph{success}, and ``abstraction'' is an interpretation licensed by
the task structure and anatomy, not a decoded representation.

\paragraph{Consolidation sits at the mesoscale of multi-step integration.}
The inference-specific $\beta$ component is multiscale and mildly
mesoscale-favouring: it cleared the matched-strength null most strongly at
intermediate diffusion scales ($\tau\approx2.6$, $p=0.014$; $\tau\approx6.8$,
$p=0.010$; leave-one-out robust), with an effect size modestly larger than at the
finest scale (observed-minus-null gap $+0.11\rightarrow+0.15$), while every
control band ($\alpha,\delta,\theta$, low-$\gamma$, high-$\gamma$) stayed null at
the mesoscale
% source: N2.2; consolidation_arc tau-sweep (audit_103d); report 2026-06-22_arc-mesoscale-inference-null
. The mesoscale is the natural scale of multi-step paths through the network, and
multi-step chaining (A$>$B, B$>$C $\Rightarrow$ A$>$D) is exactly transitive
inference. The consolidation signal sits at the scale of the computation it is
meant to reflect. (Mesoscale-favouring, not mesoscale-exclusive: the effect is
significant across scales.)

\paragraph{Anchored in the orbitofrontal map, an anchor set by learning itself.}
The encoding component concentrated in OFC, and that anchor is established by
learning rather than by reasoning. The pure memorising phase
(rest $\rightarrow$ learning $\rightarrow$ rest, with no inference demanded)
already leaves its own persistent trace, clearing the matched-strength null in
$\alpha$ ($+0.22$, $p=0.014$) and $\beta$ ($+0.25$, $p=0.014$; $\delta$ marginal,
$p=0.097$), and the $\beta$ learning-trace localises to OFC --- the same hotspot
the test-phase trace uses ($R=1000$, $q=0.010$, low-strength, non-hub)
% source: N2.6; consolidation_arc (T_learn) + localization (audit_110/111); report 2026-06-22_learning-phase-own-trace
. OFC is therefore anchored by both task phases; inference adds the $\beta$-only
component on top. Because the encoding term uses the shorter learning phase, it
carries no test/learn length asymmetry, so the duration confound below never
touched it.

\paragraph{Memory and inference ride different rhythms.}
Reading content by rhythm, $\alpha$ kept the memorised material and only that
(encoding-only), whereas $\beta$ kept the memorised material \emph{and} the
inferred relations (\tab{tab:arc}). This $\beta$-only inference result is
duration-controlled. The test phase is $1.35$--$2.53\times$ longer than the
learning phase in every patient, and $f$ inherits that asymmetry; we therefore
discarded a naive length-matched null built on a truncated test recording after
its negative control failed ($\delta$, dead-null at full length, $p\approx0.54$,
false-positived at $p\approx0.014$ under truncation)
% source: N2.5; arc_lenmatched_null (audit_103b) -- RETIRED as invalid
. In its place, a per-patient regression of the full-length inference term
against the test/learn length ratio is artifact-free: $\beta$ does not scale with
duration (Spearman $\rho\approx+0.25$, $p\approx0.49$) whereas $\alpha$ carries
the only duration-tracking trend ($\rho\approx+0.53$). Combined with the
full-length matched-strength null that $\beta$ already clears, the
inference-specific $\beta$ arc is duration-robust and $\alpha$ is
duration-suspect.

\paragraph{Whole-brain averaging conceals a focal cingulate memory trace at
low-\texorpdfstring{$\gamma$}{gamma}.}
There is no whole-brain low-$\gamma$ trace, yet an absolute within-region test
(no demeaning) found a genuine focal encoding trace in the cingulate at
low-$\gamma$ ($\rho\approx+0.39$, 8/8 patients, $q\approx0.035$; with the
seizure zone excluded $\rho\approx+0.51$ in 7/8 --- a larger effect on fewer
contacts, so not epilepsy-driven; cingulate inference was null here, so the trace
is encoding-specific)
% source: N2.4; inference_localization/within_system_trace_*.csv (audit_112)
. A real trace that cohort-wide averaging conceals is the cleanest demonstration
of why a per-region, multiscale read-out matters. Being encoding-only, it is
duration-immune.

\paragraph{Inference leans toward the cingulate.}
Feeding the localiser the inference-specific per-pair vectors, the inference
component leaned toward cingulate cortex rather than OFC --- but this is a
directional hint, not an established location. It was Benjamini--Hochberg
borderline at full length and did not survive the matched-length strength null
($p\approx0.12$--$0.18$, $q\approx0.42$--$0.76$), so its full-length significance
was duration-assisted even though its sign was consistent
% source: N2.3/N2.5; inference_localization (audit_110) + lenmatched (audit_113b)
. The inference component is duration-robust at the whole-brain level but remains
anatomically distributed with no established hotspot.

\paragraph{One cingulate region carries different content in different rhythms.}
Read together, the cingulate behaves like a hub that switches its content with
the rhythm: memory at low-$\gamma$ (solid, above) and, more tentatively,
inference at $\beta$ (a hint). We present this multiplexing as suggestive,
precisely because the inference-location half rests on a directional hint.

% ---- Table 2 ----
\begin{table}[t]
\centering
\caption{\textbf{Consolidation arc by band.} Encoding echo
($T_{\text{learn}}$) and inference-specific persistence
($T_{\text{inf}\cdot e}$, encoding partialled out), under the strength-matched
null. $\alpha$ keeps encoding only; $\beta$ keeps encoding and inference.}
\label{tab:arc}
\begin{tabular}{lcc}
\toprule
Band & Encoding echo $T_{\text{learn}}$ & Inference-specific $T_{\text{inf}\cdot e}$ \\
\midrule
$\alpha$       & $+0.22$\;($p=0.014$) & null \\
$\beta$        & $+0.25$\;($p=0.014$) & persists\;($p\approx0.007$) \\
$\delta$       & marginal\;($p=0.097$) & null \\
$\theta$       & null & null \\
low-$\gamma$   & null$^{\dagger}$ & null \\
high-$\gamma$  & null & null \\
\bottomrule
\end{tabular}
\par\smallskip
{\footnotesize $^{\dagger}$ No \emph{whole-brain} low-$\gamma$ echo, but a focal
\emph{within-region} cingulate encoding trace is present (\tab{tab:loc}).
Source: consolidation\_arc (audit\_103, \_103d); arc\_null\_per\_patient.csv.}
\end{table}

% ---- Table 3 ----
\begin{table}[t]
\centering
\caption{\textbf{Where the content concentrates}, with the duration control.
Matched-strength localiser, $R=1000$. Encoding anatomy is robust and
duration-invariant; the inference location is a directional hint that the
duration control demotes.}
\label{tab:loc}
\begin{tabular}{llll}
\toprule
Component & Region & BH $q$ ($R{=}1000$) & Duration control \\
\midrule
Encoding ($\beta$)                 & OFC       & $0.010$ & survives ($q=0.050$, epi-excl.) \\
Learning-phase trace ($\beta$)     & OFC       & $0.010$ & survives ($q=0.050$) \\
Inference ($\beta$)                & cingulate & borderline & \textbf{fails} ($q\approx0.42$--$0.76$) $\rightarrow$ hint \\
Encoding within-region (low-$\gamma$) & cingulate & $0.035$\;($\rho{\approx}0.39$, 8/8) & immune (encoding-only) \\
\bottomrule
\end{tabular}
\par\smallskip
{\footnotesize Source: localization\_atlas (audit\_83/92); inference\_localization
(audit\_110/112/113b). ANATOMY\_LEDGER.}
\end{table}

% =====================================================================
\subsection{Propagator-inspired markers of epileptogenic tissue}

The same operator that reveals the cognitive trace also speaks to the clinic.
Read not as a cross-phase distance but as inter-contact diffusion affinity, the
Laplacian propagator exposes the seizure network.

\paragraph{Epileptogenic tissue is organised relationally, not contact by
contact.}
A node-by-node ``is this contact epileptic?'' marker does not exist beyond
trivial confounds: within a patient it reduces to hubness and electrode depth,
and across patients nothing transfers. Read relationally, however, the
epileptic contacts form a strength-independent, spatially-irreducible diffusion
community --- a group that co-diffuses --- that survives strength- and
space-matched nulls. As a marker, a per-node read-out of how strongly a contact
co-diffuses with known seizure contacts recovers \emph{distant} seizure-onset
contacts (off-shaft, with spatial proximity removed by construction, since the
connectivity substrate is volume-conduction immune), in the $\delta$ band, stable
across rest and task, against a label-shuffle null at $p\approx0$
% source: N3.1; reports 2026-06-08/2026-06-12 (audit_89/99/101/102)
. (An all-contacts version was a proximity tautology and is withdrawn; the honest
metric is leave-one-shaft-out off-shaft discovery.)

\paragraph{Fused across rhythms, the propagator yields a calibrated seizure-zone
probability.}
Concatenating propagator-only features (heat affinity at multiple scales,
personalised PageRank, Katz, communicability, commute/resistance,
diffusion distance, seed-relative segregation), with no coordinates, across all
six bands, and training a leave-one-patient-out logistic model from $k=3$ seeds,
yields a calibrated probability that each contact lies in the seizure-onset zone.
Across the cohort this reached mean AUC~$\approx0.81$, precision@5~$\approx60\%$
($\sim$sevenfold over chance), with 9/10 patients above chance and a clean
label-shuffle null (\tab{tab:detector})
% source: N3.2; epi_propagator_detector/detector_lopo_per_patient.csv (audit_117)
. The lift comes from cross-band fusion rather than richer single-band features;
a gradient-boosted variant raised AUC without improving precision and was
rejected as overfit-prone at $n=10$, in favour of the interpretable logistic. (A
first draft scored $0.97$ through a label leak in the segregation feature; this
was caught and fixed, and the shuffle null now certifies the pipeline.)

\paragraph{Finding the focus without a seed remains the open frontier.}
The clinically most valuable version --- localising the seizure zone with no
known seizure contacts --- is under active investigation
% source: N3.3; audit_118/119/120 (running)
. Prior label-free work leads us to expect it to be less precise, because
selecting which co-diffusing community is the pathological one, at zero labels,
is hard; we report no seed-free verdict here.

\paragraph{Two populations: tight communities and network hubs.}
Detector performance is bimodal and must be read per patient. In a
strong-community subgroup the seizure zone is a coherent diffusion community and
detection is excellent (AUC $\approx0.90$--$0.99$, precision@5 up to $100\%$); in
a hub subgroup (Pat\_10, Pat\_15, right-hemisphere implants where the seizure
zone coincides with network hubs rather than a community) it is marginal
(\tab{tab:detector})
% source: N3.4; detector_lopo_per_patient.csv
. The cohort mean can therefore mislead, and we report the split explicitly.

\paragraph{Scope, and an honest ceiling.}
All of this is validated against clinical seizure-onset labels, not against
resection and surgical outcome (Engel/ILAE); the marker is a triage/shortlist
tool, not a from-scratch localiser. The precision ceiling ($\sim60\%$) is
structural, set by the rarity of the target times ranker quality, and no filter
rescues it --- a white-matter filter even hurts, because $\sim41\%$ of
seizure-onset contacts are white-matter-labelled. High-ranked unmarked contacts
(e.g.\ hippocampus/amygdala, medial OFC) are hypotheses, none outcome-validated.

% ---- Table 4 ----
\begin{table}[t]
\centering
\caption{\textbf{Seed-based propagator detector} (leave-one-patient-out,
$k=3$ seeds, six-band fusion, interpretable logistic). Calibrated $P(\mathrm{SOZ})$
per contact; report per patient, not the cohort mean alone.}
\label{tab:detector}
\begin{tabular}{ll}
\toprule
Metric / subgroup & Value \\
\midrule
Mean AUC (cohort)                 & $\approx0.81$ \\
Precision@5 (cohort)              & $\approx60\%$ ($\sim$7$\times$ chance) \\
Patients above chance             & 9/10 \\
Strong-community subgroup         & AUC $0.90$--$0.99$; precision@5 up to $100\%$ \\
Hub subgroup (Pat\_10, Pat\_15)   & marginal \\
\bottomrule
\end{tabular}
\par\smallskip
{\footnotesize Source: epi\_propagator\_detector/ (audit\_117);
README + detector\_lopo\_per\_patient.csv. Per-patient strip is a figure, not a
cohort scalar.}
\end{table}

% =====================================================================
\section{Discussion}

We set out to ask whether the human brain leaves an offline trace of inferred
structure, and found one with a clear scale, spectrum, anatomy and dynamics. A
transitive-inference task reorganises resting cortical connectivity in a way that
persists, peaks in $\beta$, concentrates in (without being confined to) the
orbitofrontal map, rides healthy cortex, and is held as a sustained state. The
persistence carries an inference-specific component beyond encoding, sits at the
mesoscale of multi-step integration, and is anchored in an OFC map that learning
itself establishes. We read this as an offline signature of \emph{abstraction}:
the brain keeps the relations it computed, not merely the pairs it saw.

The central methodological lesson is that this structure is only visible across
scales. It is the weakest band at the level of raw connectivity, invisible to a
connection-by-connection comparison, and --- in its $\alpha$ component --- missed
by the leading-mode spectral embeddings that spectral clustering and PCA rely on.
A multiscale diffusion read-out recovers it and, by dissociating the per-pair and
global-mode probes across bands, shows that the two views are non-redundant.
Network reorganisation that lives between scales is not a niche case; it is what a
single-scale or global-spectral analysis is structurally unable to see.

That the persistence is \emph{held} rather than bursty is, to our knowledge, a
new characterisation. We searched hard for transient, replay-like re-entry of the
task state across representations, bands and timescales down to the ripple range,
and found none; what we found instead is a sustained, non-ergodic reorganisation
the resting network settles into and holds. We are careful not to over-read this:
it is a held connectivity state, and we deliberately do not claim that rest
becomes a tighter attractor (only a weak trend), nor --- absent performance data
--- that the trace predicts inference success.

Finally, the same operator localises epileptogenic tissue, with a relational,
strength-orthogonal, proximity-immune marker that finds distant seizure-onset
contacts and, fused across bands, yields a calibrated triage probability. The
honest scope is a clinical-label, not an outcome, validation, with a structural
precision ceiling and a genuine two-population split; we present it as the
clinical face of one operator rather than a finished localiser. Across cognition
and clinic, a single multiscale diffusion read-out of intracranial connectivity
recovers structure that conventional analyses miss.

Limitations are stated throughout and summarised here: $n=10$, with irreducible
OFC coverage of five patients; one sanctioned cohort-anti patient at $\beta$
(Pat\_15, a right-hemisphere-only implant); no behavioural/performance data, by
which the abstraction reading is an interpretation and not a brain--behaviour
result; and an inference component that is duration-robust at the whole-brain
level but anatomically thin.

% =====================================================================
\section{Methods}
% (Summary scaffold; expand to journal length. Methods carries everything the
%  Results deliberately do not describe.)

\paragraph{Participants and task.}
Ten patients undergoing stereo-EEG monitoring
(\TODO{IDs, ages, implant details, ethics/consent statement, recording site}).
A transitive-inference paradigm with four phases: pre-task rest, a learning phase
in which adjacent premise pairs are shown, a test phase requiring inference on
novel non-adjacent pairs, and post-task rest
(\TODO{phase durations, trial counts, exclusion of Pat\_10 task rows 53--55}).
Performance/accuracy was not recorded and is unavailable.

\paragraph{Connectivity substrate.}
Functional connectivity was estimated as the band-averaged magnitude of the
imaginary part of coherency, $\langle|\mathrm{ImCoh}|\rangle_f$, which is immune
to zero-lag volume conduction by construction. Per-frequency-bin magnitudes are
averaged within band after the nonlinear transform (order of operations matters).
Bands: $\delta,\theta,\alpha,\beta$, low-$\gamma$, high-$\gamma$
(\TODO{exact band edges from config}). Sampling-rate differences (Pat\_03 at
1024~Hz) are absorbed at the spectral-estimation layer.

\paragraph{Laplacian diffusion hierarchy.}
From the connectivity graph we form the density operator
$\rhat=e^{-\tau\hat L}/Z$, a diffusion propagator over the graph Laplacian
$\hat L$, defining a scale-continuous renormalisation of the network indexed by
diffusion time $\tau$; the canonical scale is $\tau=1/\lambda_{\max}$. A
$\tau$-sweep confirms the trace is fine-scale and stable around this value, with
coarse-$\tau$ ``gains'' identified as a collapse artifact
% source: N1.5; audit_121
.

\paragraph{Two probes.}
(i) A per-pair cophenetic distance \rhocoph{} read from the diffusion-induced
hierarchy --- multiscale and locally interpretable. (ii) A global leading-mode
Grassmann distance \dG{} between leading-$k$ subspaces --- the view that spectral
clustering/PCA approximate. Cross-phase similarity uses
$\rho^{\mathrm{coph}}=\mathrm{Spearman}$ over the upper-triangle of the cophenetic
matrices; partition-cut metrics (ARI/NMI) are not used.

\paragraph{Cross-phase statistic and sign convention.}
Persistence is quantified with the triangle statistic
$T_d = d(\text{pre},\text{task}) - d(\text{task},\text{post})$, so that $T_d>0$
denotes a trace. One-sided Wilcoxon tests (alternative ``greater'') at the cohort
level; per-patient trace counts are $(T_d>0)$.

\paragraph{Control battery (C1--C5).}
C1, within-baseline split-half; C2, a drift floor; \textbf{C3, a strength-matched
surrogate} that preserves each node's total connectivity while destroying
higher-order structure --- the binding referee for every cohort claim (split-
baseline surrogate, audit\_63); C4, cross-probe restriction; C5,
epileptic-zone exclusion (with a size-matched node-decimation control to separate
tissue-specific from generic node-count effects). Surrogate upper-tail
$p=\mathrm{mean}(s\ge T_{\text{obs}})$, $R=1000$ for localisation. Multiplicity
within a family by Benjamini--Hochberg; leave-one-out maxima reported (no single-
patient-$p$-driven claims).

\paragraph{Localisation.}
Per-system concentration of the per-pair trace above each patient's own demeaned
baseline, across nine a-priori cortical systems, with a matched-strength null
($R=1000$), rank concordance, leave-one-out and shaft-collapse robustness checks.

\paragraph{Dynamical (held-state) analysis.}
Sliding windows (20~s; down to $\sim0.2$~s for high-$\gamma$); per-window
similarity to the task configuration relative to the pre-task baseline; tested
against a time-shuffled, node-permuted placebo null; replicated across five
Laplacian read-outs (audit\_124--141).

\paragraph{Epileptogenic marker.}
Propagator-only features per band, no coordinates; leave-one-patient-out logistic
from $k=3$ seeds; calibrated $P(\mathrm{SOZ})$; label-shuffle null; evaluated
off-shaft (leave-one-shaft-out) to avoid the proximity tautology
(\TODO{feature list, detector config from epi\_propagator\_detector/README.md}).

\paragraph{Data and code availability.}
\TODO{repository / cache locations; audit script list; CSV sources cited inline.}

% =====================================================================
\section*{Figure legends (to produce)}
% Drawn from each headline's figure ideas; build from cached CSVs, PDF vector only.
\begin{description}
\item[Fig.~1 (method + band resolution).] The diffusion hierarchy and the two
  probes; the three-layer band-resolution strip (raw FC $\rightarrow$ raw
  $D(\tau)$ $\rightarrow$ cophenetic) showing $\beta$ resolved only at the
  cophenetic layer; the band$\times$probe dissociation of \tab{tab:bandprobe}.
\item[Fig.~2 (the trace).] Cohort $\beta$ trace forest/gate plot per band, with an
  epi-excluded (X-epi) variant alongside; the held-state time-course
  (post- vs.\ pre-task rest) with the transient-negative panel. No cross-probe
  (C4) figure.
\item[Fig.~3 (anatomy + tissue).] OFC system-scale concentration on a bilateral
  brain (system, not region; ``hotspot not container'' annotation in the caption);
  gray--gray vs.\ WM--WM vs.\ seizure-core pair-class contrast.
\item[Fig.~4 (flagship: abstraction).] Four-phase schematic (encoding vs.\
  inference); the inference-specific $\beta$ persistence and its mesoscale
  signature; encoding$\rightarrow$OFC vs.\ inference$\rightarrow$cingulate
  double-dissociation brain (with the ``inference duration-downgraded'' caption);
  the low-$\gamma$ within-region cingulate ``averaging-hides-it'' panel.
\item[Fig.~5 (clinic).] Diffusion-community schematic (seeds $\rightarrow$ distant
  seizure-onset zone on a different shaft); per-patient detector strip (the
  two-population split); calibration curve and precision@$k$.
\end{description}

\end{document}
